\section{Related work}\label{sec:related}

\paragraph{Optimizers for logistic regression}
Minka~\cite{minka2003comparison} compared gradient, conjugate-gradient, quasi-Newton
and Newton (iteratively reweighted least squares) methods for logistic regression
and found Newton-type methods fastest for small and medium dimension, with the cost
of the dense Hessian becoming limiting as $d$ grows. Lin, Weng and
Keerthi~\cite{lin2008trust} proposed a trust-region Newton method in which each
subproblem is solved by CG with Hessian-vector products; it became the default
large-scale solver in LIBLINEAR and is the closest relative of the second phase of
ASH-LR. Limited-memory BFGS~\cite{liu1989limited} is the other standard
matrix-free choice and serves as a baseline here. Bach~\cite{bach2010self} showed
that the logistic loss is generalized self-concordant, which explains why Newton's
method behaves well for this problem in practice.

\paragraph{Inexact Newton methods}
Dembo, Eisenstat and Steihaug~\cite{dembo1982inexact} established that Newton's
method retains local linear, superlinear or quadratic convergence when the linear
system is solved only to a relative residual $\eta_k$ (the \emph{forcing term}),
provided $\eta_k$ is bounded away from one, tends to zero, or tends to zero as
$O(\norm{g_k})$, respectively. Eisenstat and Walker~\cite{eisenstat1996forcing}
studied adaptive choices of $\eta_k$. We use the classical choice
$\eta_k=\min(0.5,\sqrt{\norm{g_k}})$, which yields superlinear local
convergence~\cite[Thm.~7.2]{nocedal2006numerical}. Conjugate
gradients~\cite{hestenes1952methods} are the natural inner solver because the
damped Hessian is symmetric positive definite. Hessian-free methods of this type
are also used in machine learning more broadly~\cite{martens2010deep,byrd2011stochastic}.

\paragraph{Damping and regularized Newton methods}
Adding $\mu I$ to the Hessian dates back to Levenberg~\cite{levenberg1944method} and
Marquardt~\cite{marquardt1963algorithm}; adapting $\mu$ by comparing actual and
predicted decrease is the standard trust-region mechanism~\cite{nocedal2006numerical}.
Li, Fukushima, Qi and Yamashita~\cite{li2004regularized} analysed regularized
Newton methods for convex problems with $\mu_k$ proportional to a power of
$\norm{g_k}$ and showed fast local convergence even with singular Hessians.
Mishchenko~\cite{mishchenko2023regularized} proved that the simple choice
$\mu_k=\sqrt{H\norm{g_k}}$, with $H$ the Lipschitz constant of the Hessian, gives a
global $O(1/k^2)$ rate together with local superlinear convergence, without line
search, matching cubic regularization~\cite{nesterov2006cubic} at the cost of one
linear solve per iteration. Doikov, Mishchenko and
Nesterov~\cite{doikov2024super} extended this to a \emph{super-universal} method
that adapts the regularization constant and the power of the gradient norm
automatically and evaluated it on logistic regression. ASH-LR borrows the
form $\mu_k=\sqrt{\theta_k\norm{g_k}}$ from this line of work and adapts
$\theta_k$ with a ratio test; the theoretical guarantees of
\cite{mishchenko2023regularized,doikov2024super} rely on their exact update rules
and on exact linear solves, and therefore do not transfer directly to ASH-LR.

\paragraph{Numerical stability}
Evaluating the softplus function $\log(1+e^{t})$ stably is a well-known
issue~\cite{higham2002accuracy}; Wiemann, Kneib and
Hambuckers~\cite{wiemann2024softplus} discuss the softplus identity in the context
of link functions for generalized linear models. We use the equivalent
\texttt{logaddexp} formulation and a sign-dependent sigmoid.

\paragraph{Adaptive second-order methods for streaming data}
Godichon-Baggioni and Werge~\cite{godichon2025adaptive} develop inversion-free
stochastic Newton methods with $O(dN)$ cost for ill-conditioned streaming problems,
illustrating both the benefit and the expense of curvature information. Our setting
is deterministic batch optimization, but the trade-off between the cost of curvature
information and the iterations it saves is the same question we study.
